\documentclass[11pt]{article}

\usepackage[utf8]{inputenc}
\usepackage[T1]{fontenc}
\usepackage{geometry}
\usepackage{amsmath}
\usepackage{amssymb}
\usepackage{booktabs}
\usepackage{hyperref}
\usepackage{xurl}
\usepackage{xcolor}
\usepackage{listings}
\usepackage{enumitem}
\usepackage{parskip}
\usepackage{graphicx}
\graphicspath{{../images/}}

\geometry{margin=1in}

\lstset{
  language=Java,
  basicstyle=\ttfamily\small,
  keywordstyle=\color{blue},
  commentstyle=\color{gray},
  stringstyle=\color{teal},
  showstringspaces=false,
  columns=fullflexible,
  frame=single,
  framerule=0.2pt,
  breaklines=true
}

\setlist[itemize]{topsep=0.3em,itemsep=0.2em}
\setlist[enumerate]{topsep=0.3em,itemsep=0.2em}

% Simple per-section source line.
\newcommand{\notesources}[1]{\par\noindent{\small\color{gray}\textit{Sources:} \sloppy #1\par}}

% Unit-wise source strings for this subject (used by slides/notes).
% Path is relative to the lecture `latex/` folder (the compilation CWD).
% OOPS with Java (CSEG1044) — unit-wise sources
%
% These are short, student-facing reference strings intended to be used with:
%   - \slidesources{...}  (in beamer slides)
%   - \notesources{...}   (in lecture notes)
%
% Style inspiration: other subjects in My-Subjects/ use semicolon-separated sources
% and include an "accessed" date for web documentation.

\newcommand{\OOPSUnitISources}{Schildt, \textit{Java: A Beginner's Guide} (10e), McGraw-Hill (Java fundamentals + OOP basics).; Oracle Java Tutorials: Getting Started + Language Basics + Classes and Objects (accessed \today).; Java SE API docs (e.g., \texttt{String}, \texttt{Scanner}) (accessed \today).}

\newcommand{\OOPSUnitIISources}{Schildt, \textit{Java: The Complete Reference} (12e), McGraw-Hill (classes, inheritance, packages, interfaces).; Bloch, \textit{Effective Java} (3e), Addison-Wesley, 2018 (design choices; interfaces vs abstract classes).; Oracle Java Tutorials: Classes and Objects + Interfaces + Packages (accessed \today).; Java SE API docs (e.g., \texttt{Comparable}, \texttt{Runnable}) (accessed \today).}

\newcommand{\OOPSUnitIIISources}{Schildt, \textit{Java: The Complete Reference} (12e) (nested/inner classes, exceptions, threads, I/O).; Oracle Java Tutorials: Nested Classes + Exceptions + Concurrency + Basic I/O (accessed \today).; Java SE API docs (e.g., \texttt{Thread}, \texttt{AutoCloseable}) (accessed \today).}

\newcommand{\OOPSUnitIVSources}{Schildt, \textit{Java: The Complete Reference} (12e) (generics, lambdas, Swing, JDBC).; Bloch, \textit{Effective Java} (3e) (generics + lambdas).; Oracle Java Tutorials: Generics + Lambda Expressions + Swing + JDBC Basics (accessed \today).; Java SE API docs (e.g., \texttt{java.util.function}, \texttt{javax.swing}, \texttt{java.sql}) (accessed \today).}

\newcommand{\OOPSUnitVSources}{Schildt, \textit{Java: The Complete Reference} (12e) (Collections Framework + wrapper classes).; Oracle Java docs: Collections Framework overview (accessed \today).; Java SE API docs (\texttt{java.util} collections; wrapper classes in \texttt{java.lang}) (accessed \today).}

\newcommand{\OOPSUnitVISources}{Git SCM: \textit{Pro Git} (2e) (version control workflow), accessed \today.; GitHub Docs (repositories, issues, pull requests), accessed \today.; Oracle Java Tutorials: Swing + JDBC (accessed \today).; JUnit 5 User Guide (accessed \today).; Apache Maven Project: Getting Started (accessed \today).; Gradle User Manual: Getting Started (accessed \today).; UML class diagrams: UML 2.x overview/spec (accessed \today).}

\newcommand{\OOPSMiscWhyInterfacesSources}{Schildt, \textit{Java: The Complete Reference} (12e) (interfaces).; Bloch, \textit{Effective Java} (3e) (prefer interfaces where appropriate).; Oracle Java Tutorials: Interfaces (accessed \today).; Java SE API docs (e.g., \texttt{Comparable}, \texttt{Runnable}) (accessed \today).}



\title{Object Oriented Programming (CSEG1044)\\Unit 01 -- Lecture 03 Notes\\I/O + Command Line + Types + Operators}
\author{Tofik Ali}
\date{\today}

\begin{document}
\maketitle
\tableofcontents

\section{Lecture Overview}
This lecture focuses on the day-to-day basics you need to write useful Java programs:
\begin{itemize}
  \item taking input (keyboard / terminal),
  \item printing output in a readable way,
  \item using command line arguments,
  \item understanding primitive data types and operators (and common pitfalls).
\end{itemize}
\notesources{\OOPSUnitISources}

\section{Syllabus Alignment (Unit 01)}
\textbf{Unit 01:} I/O statements, command line arguments, data types, variables, operators

\section{Learning Outcomes}
After this lecture, you should be able to:
\begin{itemize}
  \item Read console input using \texttt{Scanner} and (at a basic level) \texttt{BufferedReader}.
  \item Use \texttt{System.out.println} and \texttt{System.out.printf} for formatted output.
  \item Use \texttt{String[] args} to accept command line inputs and parse simple values.
  \item Choose correct primitive types and apply operators with correct precedence and casting.
\end{itemize}

\section{Key Terms (Quick)}
\begin{itemize}
  \item \textbf{Token vs line}: token = one word/number; line = full line of text.
  \item \textbf{Formatting}: printing numbers/text with a chosen format (decimal places, alignment).
  \item \textbf{Primitive types}: built-in value types like \texttt{int}, \texttt{double}, \texttt{char}, \texttt{boolean}.
  \item \textbf{Type casting}: converting a value from one type to another (widening/narrowing).
  \item \textbf{Precedence}: order in which operators are evaluated.
\end{itemize}

\section{Detailed Notes}
\subsection{Console I/O: \texttt{Scanner} (most common for beginners)}
\texttt{Scanner} reads tokens from an input source like \texttt{System.in}.
\begin{itemize}
  \item Use \texttt{nextInt()}, \texttt{nextDouble()}, \texttt{next()} (word), \texttt{nextLine()} (whole line).
  \item \textbf{Common pitfall:} mixing \texttt{nextInt()} and \texttt{nextLine()}.\par
  After reading a number, the newline remains in the buffer; the next \texttt{nextLine()} may return an empty string.\par
  Fix: call an extra \texttt{nextLine()} after reading the number.
\end{itemize}

\subsection{Console I/O: \texttt{BufferedReader} (fast, line-oriented)}
\texttt{BufferedReader} reads full lines as \texttt{String}, and you parse manually:
\begin{itemize}
  \item Good when you want complete control or faster I/O.
  \item Slightly more code than \texttt{Scanner} (because you must parse numbers using \texttt{Integer.parseInt}, etc.).
\end{itemize}

\subsection{Output formatting}
Use:
\begin{itemize}
  \item \texttt{System.out.println(...)} for simple output.
  \item \texttt{System.out.printf("x = \%d, y = \%.2f\%n", x, y);} for formatted output.
\end{itemize}
Common format specifiers:
\begin{itemize}
  \item \texttt{\%d} integer,\;\; \texttt{\%f} floating-point,\;\; \texttt{\%s} string,\;\; \texttt{\%n} newline.
\end{itemize}

\subsection{Command line arguments}
When you run a Java program, anything after the class name becomes a string in \texttt{args}.
\begin{itemize}
  \item Example: \texttt{java Sum 10 20 30}
  \item Then: \texttt{args[0] = "10"}, \texttt{args[1] = "20"}, etc.
\end{itemize}
You often parse numbers using:
\begin{itemize}
  \item \texttt{Integer.parseInt(args[i])}
  \item \texttt{Double.parseDouble(args[i])}
\end{itemize}
\textbf{Pitfall:} if input is not numeric, parsing throws \texttt{NumberFormatException}.

\subsection{Data types, operators, precedence, casting}
\subsubsection{Primitive data types (quick)}
\begin{itemize}
  \item Integer types: \texttt{byte}, \texttt{short}, \texttt{int}, \texttt{long}
  \item Floating types: \texttt{float}, \texttt{double}
  \item Other: \texttt{char} (Unicode), \texttt{boolean}
\end{itemize}

\subsubsection{Operator pitfalls you must know}
\begin{itemize}
  \item \textbf{Integer division:} \texttt{5/2 == 2} (not 2.5). Use \texttt{5/2.0} or cast: \texttt{(double)5/2}.
  \item \textbf{Overflow:} \texttt{int} can overflow silently.\par
  Example: \texttt{int x = 2\_000\_000\_000; x = x + 1\_000\_000\_000;} (wraps around).
  \item \textbf{Short-circuit logic:} \texttt{\&\&} and \texttt{||} may skip evaluating the right side.\par
  This is useful for null checks: \texttt{s != null \&\& s.length() > 0}.
  \item \textbf{Precedence:} multiplication happens before addition. Use parentheses when teaching/learning.
\end{itemize}

\subsubsection{Casting}
\textbf{Widening} (safe) happens automatically: \texttt{int} $\rightarrow$ \texttt{double}.
\textbf{Narrowing} may lose data and requires explicit cast: \texttt{double} $\rightarrow$ \texttt{int}.

\section{Worked Example: Args + Scanner + Formatting}
\begin{lstlisting}
import java.util.Scanner;

public class InputDemo {
    public static void main(String[] args) {
        int maxMarks = 100;
        if (args.length > 0) maxMarks = Integer.parseInt(args[0]);

        Scanner sc = new Scanner(System.in);
        System.out.print("Enter your name: ");
        String name = sc.nextLine();

        System.out.print("Enter marks (0-" + maxMarks + "): ");
        int marks = sc.nextInt();

        double percent = marks * 100.0 / maxMarks;
        System.out.printf("Hello %s, marks=%d/%d, percent=%.2f%%%n",
                name, marks, maxMarks, percent);
        sc.close();
    }
}
\end{lstlisting}

\textbf{What to observe:}
\begin{itemize}
  \item How \texttt{args} is always a string array and needs parsing.
  \item How we used a default value (100) and optionally override it using \texttt{args[0]}.
  \item How formatted output (\texttt{printf}) improves readability.
  \item Why we used \texttt{100.0} in the formula (to avoid integer division mistakes).
\end{itemize}
\notesources{\OOPSUnitISources}

\section{In-Class Exercises (with brief solutions)}
\begin{enumerate}
  \item What is the output of \texttt{System.out.println(5/2);}?\; \textbf{Answer:} \texttt{2}.
  \item How to get \texttt{2.5} from 5 and 2?\; \textbf{Answer:} \texttt{5/2.0} or \texttt{(double)5/2}.
  \item If you run \texttt{java Demo A B}, what is \texttt{args.length}?\; \textbf{Answer:} 2.
\end{enumerate}

\section{Self-Study Practice (no solutions)}
\begin{enumerate}
  \item Write a program \texttt{SumArgs} that sums integer command line arguments.
  \item Write a program that reads a line and prints: (a) length, (b) first character, (c) uppercase version.
  \item Create 5 expressions and predict results using precedence; then verify by running Java.
\end{enumerate}

\section{Checkpoint Questions (with sample answers)}
\textbf{Q1.} Why is \texttt{Scanner.nextLine()} tricky after \texttt{nextInt()}?\par
\textbf{Sample answer:} \texttt{nextInt()} leaves the newline in the input buffer, so the next \texttt{nextLine()} may read the leftover newline and return an empty string. We can consume the newline by calling \texttt{nextLine()} once after reading the number.

\vspace{0.6em}
\textbf{Q2.} Why is \texttt{5/2} not 2.5 in Java?\par
\textbf{Sample answer:} Both operands are integers, so Java performs integer division and discards the fractional part. If either operand is \texttt{double}, it performs floating-point division.

\end{document}
